%
% Template Laporan Skripsi/Thesis
%
% @author  Andreas Febrian, Lia Sadita
% @version 1.03
%
% Dokumen ini dibuat berdasarkan standar IEEE dalam membuat class untuk
% LaTeX dan konfigurasi LaTeX yang digunakan Fahrurrozi Rahman ketika
% membuat laporan skripsi. Konfigurasi yang lama telah disesuaikan dengan
% aturan penulisan thesis yang dikeluarkan UI pada tahun 2008.
%

%
% Tipe dokumen adalah report dengan satu kolom.
%
\documentclass[12pt, a4paper, onecolumn, oneside, final]{report}
%\raggedbottom

% Load konfigurasi LaTeX untuk tipe laporan thesis
\usepackage{etc/uithesis}

% Daftar pemenggalan suku kata dan istilah dalam LaTeX
\input{etc/hype.indonesia}

% Load konfigurasi khusus untuk laporan yang sedang dibuat
%-----------------------------------------------------------------------------%
% Informasi Mengenai Dokumen
%-----------------------------------------------------------------------------%
%
% Judul laporan.
\var{\judul}{Judul Karya Ilmiah Anda}
%
% Tulis kembali judul laporan, kali ini akan diubah menjadi huruf kapital
\Var{\Judul}{Judul Karya Ilmiah Anda}
%
% Tulis kembali judul laporan namun dengan bahasa Ingris
\var{\judulInggris}{Your Scientific Publication Title}

%
% Tipe laporan, dapat berisi Skripsi, Tugas Akhir, Thesis, atau Disertasi
\var{\type}{skripsi}
%
% Tulis kembali tipe laporan, kali ini akan diubah menjadi huruf kapital
\var{\Type}{Skripsi}
%
% Tulis nama penulis
\var{\penulis}{Nama Lengkap Anda}
%
% Tulis kembali nama penulis, kali ini akan diubah menjadi huruf kapital
\Var{\Penulis}{Nama Lengkap Anda}
%
% Tulis NPM / NIM penulis
\var{\npm}{Nomor Anda}
\var{\nim}{Nomor Anda}
%
% Tuliskan Fakultas dimana penulis berada
\Var{\Fakultas}{Fakultas Anda}
\var{\fakultas}{Fakultas Anda}
\var{\websiteFakultas}{www.abc.uingusdur.ac.id}
\var{\emailFakultas}{abc@uingusdur.ac.id}
%
% Tuliskan Program Studi yang diambil penulis
\Var{\Program}{Jurusan Anda}
\var{\program}{Jurusan Anda}
%
% Tuliskan tahun publikasi laporan
\Var{\bulanTahun}{Bulan Tahun}
\Var{\tahun}{2022}
%
% Tuliskan gelar yang akan diperoleh dengan menyerahkan laporan ini
\var{\gelar}{Gelar Jurusan Anda}
%
% Tuliskan tanggal pengesahan laporan, waktu dimana laporan diserahkan ke
% penguji/sekretariat
\var{\tanggalSiapSidang}{Tanggal Bulan Tahun}
%
% Tuliskan tanggal keputusan sidang dikeluarkan dan penulis dinyatakan
% lulus/tidak lulus
\var{\tanggalLulus}{Tanggal Bulan Tahun}
\var{\hariLulus}{Hari}
% Tuliskan tanggal pengesahan laporan final, waktu dimana laporan
% diserahkan ke perpustakaan
\var{\tanggalFinal}{Tanggal Bulan Tahun}
% sama dengan tanggal final
\var{\tanggalPengesahan}{Tanggal Bulan Tahun}
%
% Tuliskan pembimbing
\var{\pembimbingSatu}{Pembimbing Pertama Anda}
\var{\nipPembimbingSatu}{1234567890}
\var{\pembimbingDua}{Pembimbing Kedua Anda}
%
% Tuliskan penguji
\var{\pengujiSatu}{Penguji Pertama Anda}
\var{\nipPengujiSatu}{1234567890}
\var{\pengujiDua}{Penguji Kedua Anda}
\var{\nipPengujiDua}{1234567890}

%-----------------------------------------------------------------------------%
% Judul Setiap Bab
%-----------------------------------------------------------------------------%
%
% Berikut ada judul-judul setiap bab.
% Silahkan diubah sesuai dengan kebutuhan.
%
\Var{\kataPengantar}{Kata Pengantar}
\Var{\babSatu}{Pendahuluan}
\Var{\babDua}{Kerangka Berpikir}
\Var{\babTiga}{Judul Bab 3}
\Var{\babEmpat}{Judul Bab 4}
\Var{\babLima}{Judul Bab 5}
\Var{\kesimpulan}{Penutup}


% Daftar istilah yang mungkin perlu ditandai
\input{src/istilah}

% Awal bagian penulisan laporan
\begin{document}

%
% Sampul Laporan
%
% Sampul Laporan

%
% @author  unknown
% @version 1.10
% @edit by Andreas Febrian
% @modified by Slamet B. Aryo
%

\begin{titlepage}
    \begin{center}

        \vspace*{1.0cm}
        % judul thesis harus dalam 14pt Times New Roman
        \bo{\Judul} \\[1.0cm]

        \vspace*{0.5cm}
        % harus dalam 14pt Times New Roman
        \bo{\MakeUppercase{\Type}} \\

        \vspace*{1cm}
        % keterangan prasyarat
        Diajukan untuk memenuhi sebagian syarat memperoleh gelar \\
        \gelar\\

        \vspace*{1cm}
        \begin{figure}
            \begin{center}
                \includegraphics[width=4cm]{assets/pics/logo-color.ps}
            \end{center}
        \end{figure}

        \vspace*{1cm}
        % penulis dan nim
        Oleh:\\
        \underline{\bo{\Penulis}} \\
        \bo{NIM: \nim} \\

        %\vspace*{1.5cm}
        \vfill

        % informasi mengenai fakultas dan program studi
        \bo{
            JURUSAN \Program \\
            FAKULTAS \Fakultas\\
            UNIVERSITAS ISLAM NEGERI K.H. ABDURRAHMAN WAHID\\
            \tahun
        }
    \end{center}
\end{titlepage}


%\forceclearchapter

%
% Gunakan penomeran romawi
\pagenumbering{roman}
%\pagestyle{first-pages}

%
% load halaman judul dalam
\addChapter{HALAMAN JUDUL}
%
% Halaman Judul Laporan
%
% @author  unknown
% @version 1.10
% @edit by Andreas Febrian
% @modified by Slamet B. Aryo
%

\thispagestyle{plain}

\begin{center}
 \vspace*{0cm}

    \vspace*{1.0cm}
    % judul thesis harus dalam 14pt Times New Roman
    \bo{\Judul} \\[1.0cm]

    \vspace*{0.5cm}
    % harus dalam 14pt Times New Roman
    \bo{\MakeUppercase{\Type}} \\

    \vspace*{1cm}
    % keterangan prasyarat
        Diajukan untuk memenuhi sebagian syarat memperoleh gelar \\
    \gelar\\

    \vspace*{1cm}
        \begin{figure}
        \begin{center}
            \includegraphics[width=4cm]{assets/pics/logo-aswad.ps}
        \end{center}
    \end{figure}

    \vspace*{1cm}
    % penulis dan nim
    Oleh:\\
    \underline{\bo{\Penulis}} \\
    \bo{NIM: \nim} \\

    \vfill

    % informasi mengenai fakultas dan program studi
    \bo{
    	JURUSAN \Program \\
    	FAKULTAS \Fakultas\\
    	UNIVERSITAS ISLAM NEGERI K.H. ABDURRAHMAN WAHID\\
    	\tahun
    }
\end{center}


%\forceclearchapter

%
% setelah bagian ini, halaman dihitung sebagai halaman ke 2
%\setcounter{page}{2}
\doublespacing

%
% load halaman orisinalitas
\addChapter{LEMBAR PERNYATAAN KEASLIAN}
%
% Halaman Orisinalitas
%
% @author  Andreas Febrian
% @version 1.10
% @modified by Slamet B. Aryo
%

\chapter*{Surat Pernyataan Keaslian Karya}
\vspace*{2em}

\noindent Yang bertanda tangan di bawah ini:\\

\noindent
\begin{tabularx}{\linewidth}{@{}l c @{\hspace{.5em}}X@{}}
    Nama & : & \penulis \\
    NIM & : & \nim \\
    Judul Skripsi & : & \bo{\judul} \\
\end{tabularx}

\vspace{2em}
\noindent
menyatakan dengan sesungguhnya bahwa \MakeLowercase{\type} ini adalah benar-benar
hasil karya penulis, kecuali dalam bentuk kutipan yang telah
penulis sebutkan sumbernya.
Demikian pernyataan ini penulis buat dengan sebenar-benarnya.\\

\begin{ttdkanan}
Pekalongan, \tanggalPengesahan\\
Yang menyatakan,\\
\end{ttdkanan}

%\vspace*{1em}
\begin{adjustwidth}{7cm}{}
\fbox{
    \parbox{1.5cm}{
    Meterai\\
    10.000
    }
}
\end{adjustwidth}

\vspace*{1em}
\begin{ttdkanan}
\penulis\\
\end{ttdkanan}

\newpage

%\forceclearchapter
%
% load halaman nota pembimbing
\addChapter{LEMBAR NOTA PEMBIMBING}
%
% Halaman Nota pembimbing
%
% @author  Slamet B. Aryo
% @version 1.10
%

\chapter*{NOTA PEMBIMBING}

\vspace*{0.2cm}

\begin{spacing}{1}
\noindent
\bo{\pembimbingSatu}\\
Alamat dosen pembimbing\\
\underline{Alamat baris kedua jika ada}\\

\noindent
\begin{tabbing}
Lamp. \hspace{1em}  \= : 2 (dua) eksemplar\\
Hal 				\> : Naskah Skripsi Sdr. \penulis
\end{tabbing}

\noindent
Kepada Yth.\\
Dekan Fakultas \fakultas\\
\textit{c.q.} Ketua Jurusan \program\\
\underline{PEKALONGAN}\\[1em]
\end{spacing}

\begin{spacing}{1.215}

\noindent
\emph{Assalamualaikum Wr. Wb.}\\

\noindent
Setelah diadakan penelitian dan perbaikan seperlunya,
maka bersama ini saya kirimkan naskah \text{\type} saudara:\\

\noindent
\begin{tabularx}{\linewidth}{@{}l c @{\hspace{.5em}}X@{}}
	Nama &: & \bo{\penulis} \\
	NIM &: & \bo{\nim} \\
	Judul \Type &: & \bo{\judul} \\
\end{tabularx} \\

\vspace*{1em}

\noindent
Naskah tersebut sudah memenuhi persyaratan untuk dapat segera dimunaqosahkan.
Demikian nota pembimbing ini dibuat untuk digunakan sebagaimana mestinya.
Atas perhatiannya, saya sampaikan terima kasih.\\[0.3cm]

\noindent
\emph{Wassalamualaikum Wr. Wb.}

\begin{ttdkanan}
Pekalongan, \tanggalPengesahan \\
Pembimbing,\\[3em]

\noindent
\underline{\bo{\pembimbingSatu}}\\[0.1cm]
NIP. \nipPembimbingSatu
\end{ttdkanan}

\end{spacing}

\newpage

%\forceclearchapter
%
\addChapter{LEMBAR PENGESAHAN}
%
% Halaman Pengesahan Sidang
%
% @author  Andreas Febrian, Andre Tampubolon
% @version 1.02
%

\thispagestyle{kopkampus}
\enlargethispage{-3\baselineskip}
%\chapter*{}

% \begin{tabular}[t]{@{} l @{}}
%   \includegraphics[height=3\normalbaselineskip]{assets/pics/logo-color.png} \\
% \end{tabular}%
% \hfill
% \begin{tabular}[t]{@{} r @{}}
%   \bfseries My Name \\
%   \itshape Dept.\ of Science \\
%   \itshape Prestigious University \\
%   \itshape 000-000-0000 \\
%   \itshape email@university.edu \\
%   \itshape www.google.com \\
% \end{tabular}

% kop terinspirasi dari
% https://www.overleaf.com/latex/templates/letterhead/zkvhmkzskgdm

% \begin{figure}[!ht]
% %\begin{minipage}{\textwidth}
% \begin{center}
% \begin{spacing}{.85}
% \begin{tabular}[t]{l c}
% \multirow{5}{*}{\includegraphics[width=0.15\linewidth]{assets/pics/logo-color.png}} &
% \bo{KEMENTERIAN AGAMA REPUBLIK INDONESIA}  \\
% & \small{\bo{UNIVERSITAS ISLAM NEGERI}} \\
% & \small{\bo{K.H. ABDURRAHMAN WAHID PEKALONGAN}} \\
% & \small{\bo{FAKULTAS \Fakultas}} \\
% & \footnotesize{Jalan Pahlawan No. 52 Rowolaku Kajen Kab. Pekalongan} \\
% & \footnotesize{Website: \websiteFakultas\ Email: \emailFakultas} \\
% \end{tabular}
%
% % gambar garis kop
% \hrule width \hsize \kern 0.5mm \hrule width \hsize height 2pt
% \end{spacing}
% \end{center}
% %\end{minipage}
% \end{figure}

\chapter*{P E N G E S A H A N}

\vspace*{0.4cm}

Dekan Fakultas \fakultas\ Universitas Islam Negeri (UIN) K.H. Abdurrahman Wahid Pekalongan mengesahkan \type\ Saudara:\\

\noindent
\begin{tabularx}{\linewidth}{@{}l l @{\hspace{.5em}}X@{}}
	%\type~ini diajukan oleh&: & \\
	Nama&: & \bo{\penulis} \\
	NIM&: & \bo{\nim} \\
	%Jurusan&: & \program \\
	Judul \Type&: & \bo{\judul} \\
\end{tabularx} \\

\vspace*{1.5em}

Telah diujikan pada hari \hariLulus\ tanggal \tanggalLulus\ dan dinyatakan
\bo{\underline{LULUS}} serta diterima sebagai salah satu syarat guna
memperoleh gelar \gelar.\\[0.2cm]

\begin{center}
	Dewan Penguji,
\end{center}

\vspace*{0cm}

% \begin{tabular}{l l l l }
% 	& & & \\
% 	Pembimbing 1&: & \pembimbingSatu & (\hspace*{3.0cm}) \\
% 	& & & \\
% 	Pembimbing 2&: & \pembimbingDua & (\hspace*{3.0cm}) \\
% 	& & & \\
% 	Penguji 1&: & \pengujiSatu & (\hspace*{3.0cm}) \\
% 	& & & \\
% 	Penguji 2&: & \pengujiDua & (\hspace*{3.0cm}) \\
% \end{tabular}\\

% \vspace*{2.0cm}
%
% \begin{tabular}{ll l}
% 	Ditetapkan di&: & Depok\\
% 	Tanggal&: & \tanggalLulus \\
% \end{tabular}

\begin{center}
\begin{multicols}{2}
Penguji I\\[3em]
\underline{\bo{\pengujiSatu}}\\[-0.6em]
NIP \nipPengujiSatu

Penguji II\\[3em]
\underline{\bo{\pengujiDua}}\\[-0.6em]
NIP \nipPengujiDua
\end{multicols}

Pekalongan, \tanggalLulus\\[-0.6em]
Dekan Fakultas \fakultas\\[3em]

\underline{\bo{(Nama Dekan)}}\\[-0.6em]
NIP 123456789
\end{center}


\newpage

%\forceclearchapter
%
%
\addChapter{LEMBAR MOTTO dan PERSEMBAHAN}
\chapter*{MOTTO}

``Motto''\\

\newpage

\chapter*{PERSEMBAHAN}

Karya ini kupersembahkan kepada...\\

\newpage

%\forceclearchapter
%
%
\addChapter{ABSTRAK}
\input{src/00-frontmatter/06-abstrak}
%
%
%
% Halaman Abstrak
%
% @author  Andreas Febrian
% @author  Slamet B. Aryo
% @version 1.01
%

\chapter*{Abstrak}
%\singlespacing

\vspace*{0.2cm}

% \noindent \begin{tabular}{l l p{10cm}}
% 	Nama&: & \penulis \\
% 	Program Studi&: & \program \\
% 	Judul&: & \judul \\
% \end{tabular} \\

%\vspace*{0.5cm}
\noindent \bo{\Penulis. \judul.}

%\paragraph*{}
\lipsum[10-12]

\vspace*{0.2cm}

\noindent \bo{Kata kunci:} Kata kunci satu, kata kunci dua \\

%\onehalfspacing
\newpage

%-----------------
\chapter*{ABSTRACT}
%\singlespacing

\vspace*{0.2cm}

% \noindent \begin{tabular}{l l p{11.0cm}}
% 	Name&: & \penulis \\
% 	Program&: & \program \\
% 	Title&: & \judulInggris \\
% \end{tabular} \\

%\vspace*{0.5cm}

\noindent \bo{\Penulis. \judulInggris.}

%\paragraph*{}
\lipsum[67-69] \\

\vspace*{0.2cm}

\noindent \bo{Key words:} Keyword one, keyword two \\

%\onehalfspacing
\newpage

%
\addChapter{\kataPengantar}
%-----------------------------------------------------------------------------%
\chapter*{\kataPengantar}
%-----------------------------------------------------------------------------%

Puji syukur penulis ucapkan kepada Tuhan Yang Maha Esa, karena atas berkat
dan rahmat-Nya dapat menyelesaikan \type~ini.
Penulisan \type~ini dilakukan dalam rangka memenuhi salah satu syarat untuk
mencapai gelar \gelar~Jurusan \program~pada Fakultas \fakultas~UIN Gusdur.
Penulis menyadari bahwa tanpa bantuan dan bimbingan dari berbagai pihak,
dari masa perkuliahan sampai pada penyusunan \type~ini, sangatlah sulit bagi
penulis untuk menyelesaikan \type~ini.
Oleh karena itu, penulis mengucapkan terima kasih kepada:
\begin{enumerate}[topsep=0pt,itemsep=-1ex,partopsep=1ex,parsep=1ex]
	\item Dr.\ldots selaku Rektor UIN Gusdur.
	\item Dr.\ldots selaku Dekan \fakultas.
	\item Dr.\ldots selaku Wakil Dekan bidan Akademik dan Kelembagaan \fakultas.
	\item Dr.\ldots selaku Ketua Jurusan \program.
	\item Dr.\ldots selaku Sekretaris Jurusan \program.
	\item Dr.\ldots selaku dosen pembimbing yang telah menyediakan waktu, tenaga dan pikiran
	untuk mengarahkan penulis dalam penyusunan \type~ini.
	\item Dr.\ldots selaku Dosen Penasihat Akademik (DPA).
	\item Dr.\ldots selaku dosen penguji.
	\item Pihak X\ldots yang telah banyak membantu dalam memperoleh data yang
	diperlukan oleh penulis.
	\item Orang tua dan keluarga penulis yang telah memberikan bantuan dan
	dukungan material dan moral.
	\item Sahabat yang telah banyak membantu penulis dalam menyelesaikan \type~ini.
\end{enumerate}

Akhir kata, penulis berharap Tuhan Yang Maha Esa membalas segala kebaikan semua
pihak yang telah membantu. Semoga \type~ini membawa manfaat bagi perkembangan
ilmu.\\

\begin{ttdkanan}
 Tempat, tanggal, bulan, tahun\\[3em]

 \noindent\penulis
\end{ttdkanan}


% \vspace*{0.1cm}
% \begin{flushright}
% Depok, \tanggalSiapSidang\\[0.1cm]
% \vspace*{1cm}
% \penulis
%
% \end{flushright}

%
%
% Daftar isi, gambar, tabel, dan kode
%
\phantomsection %hack to make them clickable
%\singlespacing
\tableofcontents
%\onehalfspacing
\clearpage
\phantomsection %hack to make them clickable
%\singlespacing
\listoffigures
%\onehalfspacing
\clearpage
\phantomsection %hack to make them clickable
%\singlespacing
\listoftables
%\onehalfspacing
\clearpage
%\phantomsection %hack to make them clickable
%\addcontentsline{toc}{chapter}{\lstlistlistingname}
%\singlespacing
%\lstlistoflistings
%\listoflistings
%\onehalfspacing
%\clearpage
% Jika penomoran romawi selesai di ganjil
%\naiveoddclearchapter
% Jika penomoran romawi selesai di genap
%\naiveevenclearchapter

%
% Gunakan penomeran Arab (1, 2, 3, ...) setelah bagian ini.
%
\pagenumbering{arabic}
%\pagestyle{standard}
% \setlength{\belowcaptionskip}{+2pt}


%\setoddevenheader
%-----------------------------------------------------------------------------%
\chapter{\babSatu}
%-----------------------------------------------------------------------------%
Pada bab ini, akan dijelaskan tentang latar belakang dan permasalahan yang diselesaikan pada penelitian ini.


%-----------------------------------------------------------------------------%
\section{Latar Belakang}
%-----------------------------------------------------------------------------%
Tentukan latar belakang dari penelitian Anda di sini (\f{background}).
Karena kita berkuliah di perguruan tinggi Islam, maka sepantasnya ada minimal
satu dalil yang dimasukkan, entah itu ayat Al-Quran ataupun \d{h}adi\ardot{s}.
Contohnya sebagai berikut:

\setRTL
\noindent \textarabic{\quranayah*[5][2]} %cukup memasukkan nomor surat dan ayat
\setLTR
% \end{arabic}
\vspace*{12pt}

\begin{spacing}{1.125}
``Wahai orang-orang yang beriman, janganlah kamu melanggar syiar-syiar (kesucian) Allah,
jangan (melanggar kehormatan) bulan-bulan haram, jangan (mengganggu) \textit{hadyu} (hewan-hewan kurban)
dan \textit{qal\=a'id} (hewan-hewan kurban yang diberi tanda), dan jangan (pula mengganggu)
para pengunjung Baitulharam sedangkan mereka mencari karunia dan rida Tuhannya! Apabila kamu
telah bertahalul (menyelesaikan ihram), berburulah (jika mau). Janganlah sekali-kali kebencian(-mu)
kepada suatu kaum, karena mereka menghalang-halangimu dari Masjidilharam, mendorongmu berbuat
melampaui batas (kepada mereka). Tolong-menolonglah kamu dalam (mengerjakan) kebajikan dan takwa,
dan jangan tolong-menolong dalam berbuat dosa dan permusuhan.
Bertakwalah kepada Allah, sesungguhnya Allah sangat berat siksaan-Nya.''
(Al-M\={a}'idah [5]: 2)\par
\end{spacing}
\vspace*{6pt}

\lipsum[44]

%-----------------------------------------------------------------------------%
\section{Permasalahan}
%-----------------------------------------------------------------------------%
Sebutkan permasalahan penelitian Anda dari latar belakang tersebut.

%-----------------------------------------------------------------------------%
\subsection{Rumusan Masalah}
%-----------------------------------------------------------------------------%
Berikut ini adalah rumusan permasalahan dari penelitian yang dilakukan:
\begin{enumerate}
\item Bagaimana cara membuat pertanyaan penelitian?

\item Seperti apa pertanyaan yang bagus?

\item Bagimana cara menentukan pertanyaan yang tepat?
\end{enumerate}


%-----------------------------------------------------------------------------%
\subsection{Batasan Permasalahan}
%-----------------------------------------------------------------------------%
Berikut ini adalah asumsi yang digunakan sebagai batasan penelitian ini:
\begin{enumerate}
	\item Salah satu batasannya adalah, ini hanya \f{template}.
\end{enumerate}


%-----------------------------------------------------------------------------%
\section{Tujuan Penelitian}
%-----------------------------------------------------------------------------%
Berikut ini adalah tujuan penelitian yang dilakukan:
\begin{enumerate}
	\item Untuk memberikan \f{template} yang dapat mempermudah skripsi orang lain.
\end{enumerate}


%-----------------------------------------------------------------------------%
\section{Posisi Penelitian}
%-----------------------------------------------------------------------------%
Sebutkan posisi penelitian Anda.

Pertama, skripsi \citefirstlastauthor{Hani2022} \parencite*{Hani2022} berjudul
``\citefield{Hani2022}{title}'' yang menganalisis\ldots
\lipsum[11]

\begin{figure}
	\centering
	\caption{Penjelasan singkat terkait gambar.}
	\includegraphics[width=\textwidth]{assets/pics/makara.png}\\
	{\footnotesize Sumber: \textcite{Sugiyono2008}\par}
	\label{fig:research_position}
\end{figure}

Jelaskan \pic~\ref{fig:research_position} di sini.


%-----------------------------------------------------------------------------%
\section{Langkah Penelitian}
%-----------------------------------------------------------------------------%
Berikut ini adalah langkah penelitian yang telah dilakukan:
\begin{enumerate}
	\item Tinjauan literatur \\
	Pada tahap ini, dipelajari teori-teori yang terkait dengan penelitian ini untuk mendapatkan konsep dasar yang dibutuhkan dalam mencapai tujuan penelitian.
	\item Analisis implementasi dan kesimpulan \\
	Pada tahap ini, digunakan studi kasus untuk analisis terkait kegunaan \f{template}. Setelah melakukan analisis tersebut, ditarik kesimpulan keseluruhan dari penelitian ini.
\end{enumerate}


%-----------------------------------------------------------------------------%
\section{Sistematika Penulisan}
%-----------------------------------------------------------------------------%
Sistematika penulisan laporan adalah sebagai berikut:
\begin{description}[style=unboxed,leftmargin=2.5cm]
% \begin{itemize}
	\item [Bab 1:] \babSatu \\
	    Bab ini mencakup latar belakang, cakupan penelitian, dan pendefinisian masalah.
	\item [Bab 2:] \babDua \\
	    Bab ini mencakup pemaparan terminologi dan teori yang terkait dengan penelitian berdasarkan hasil tinjauan pustaka yang telah digunakan, sekaligus memperlihatkan kaitan teori dengan penelitian.
	\item [Bab 3:] \babTiga \\
	    Apa itu Bab 3?
	\item [Bab 4:] \babEmpat \\
		Apa itu Bab 4?
	\item [Bab 5:] \babLima, jika ada \\
	    Apa itu Bab 5?
	\item [Bab 6:] \kesimpulan \\
	    Bab ini mencakup kesimpulan akhir penelitian dan saran untuk pengembangan berikutnya.
% \end{itemize}
\end{description}

%\clearchapter
%-----------------------------------------------------------------------------%
\chapter{\babDua}
%-----------------------------------------------------------------------------%
Untuk memulai penelitian, dibutuhkan kerangka berpikir yang sesuai untuk permasalahan yang ingin dipecahkan. Untuk membentuk kerangka berpikir yang sesuai, perlu dikaitkan dengan hasil studi literatur yang telah dilakukan. Oleh karena itu, pada bab ini, akan dijelaskan hasil studi literatur yang telah dilakukan yang telah dikaitan dengan kerangka kerja untuk penelitian ini.


%-----------------------------------------------------------------------------%
\section{\f{Sample Theory}}
%-----------------------------------------------------------------------------%
Teori yang saya gunakan dalam penelitian ini adalah teori dari \textcite{Williams2002organizational}.

\begin{figure}
	\centering
	\caption{Pendukung teori yang dipakai}
	\includegraphics[width=0.6\textwidth]{assets/pics/makara.png}\\
	{\footnotesize Sumber: \textcite{Sugiyono2008}\par}
	\label{fig:sample}
\end{figure}

Bahwa \lipsum[34]

Jelaskan \pic~\ref{fig:sample}.

Oleh karena itu, \lipsum[35-37]

%-----------------------------------------------------------------------------%
\section{Keterkaitan Teori Dengan Penelitian}
%-----------------------------------------------------------------------------%
\begin{figure}
	\centering
	\caption{Keterkaitan konsep hasil studi literatur terhadap penelitian}
	\includegraphics[width=0.8\textwidth]{assets/pics/makara.png}\\
	\label{fig:research_concept_map}
	{\footnotesize Sumber: \textcite{Sugiyono2008}\par}
\end{figure}

Jelaskan \pic~\ref{fig:research_concept_map} di sini.

%\clearchapter
%-----------------------------------------------------------------------------%
\chapter{\babTiga}
%-----------------------------------------------------------------------------%
Apa itu Bab 3?


%-----------------------------------------------------------------------------%
\section{Contoh}
%-----------------------------------------------------------------------------%
\begin{enumerate}
	\item Tahap 1. \\ %\lstinputlisting[language=Java, caption=Contoh kode, label=code:sample]{assets/3-sample.java}
	Jelaskan teori menurut rujukan contoh \citep*{sahroni2018fikih} dengan semua nama penulis.

	\lipsum[6-8]

	\item Tahap 2.
\end{enumerate}

Tulis kesimpulan di sini. Contoh rujukan sitasi \citep{UINJAKARTARoikhan2017} yang akan
ditampilkan pada halaman daftar pustaka.

Kembali merujuk \citep{sahroni2018fikih} tetapi dengan nama penulis yang digabung.

%\clearchapter
%-----------------------------------------------------------------------------%
\chapter{\babEmpat}
%-----------------------------------------------------------------------------%
Apa itu Bab 4?


%-----------------------------------------------------------------------------%
\section{Contoh}
%-----------------------------------------------------------------------------%
\begin{enumerate}
	\item Tahap 1. \\ %\lstinputlisting[language=Python, caption=Contoh kode, label=code:sample2]{assets/4-sample.py}
	Jelaskan %\lst~\ref{code:sample2} di sini.

	\item Tahap 2.
\end{enumerate}

Tulis kesimpulan di sini.

\begin{kutipan}
 \lipsum[22]
 \ldots (\narsumSatu, 2022)
\end{kutipan}

Penjelasan hasil wawancara dengan \narsumSatu~ \lipsum[23]

%\clearchapter
%-----------------------------------------------------------------------------%
\chapter{\babLima}
%-----------------------------------------------------------------------------%
Apa itu Bab 5?


%-----------------------------------------------------------------------------%
\section{Contoh}
%-----------------------------------------------------------------------------%
\begin{enumerate}
	\item Nomor 1.
	\item Nomor 2.
\end{enumerate}

\begin{table}
    \onehalfspacing
    \caption{Data informan pada \lokus}
	\begin{center}
	\begin{tabularx}{\textwidth}{|c|X|X|}
    \hline
    \bo{No.} & \bo{Nama} & \bo{Peran} \\ [0.5ex]
    \hline
    1 & \narsumSatu & Pimpinan \lokus~ \\
    \hline
    2 & \narsumDua & Anggota pada \lokus~dan anggota masyarakat \\
    \hline
    3 & \narsumTiga & Pengawas \lokus~ \\ [1ex]
    \hline
    \multicolumn{3}{l}{\footnotesize\textit{Sumber: Catatan penulis}} \\
    \end{tabularx}
    \end{center}
	\label{table:sample}
\end{table}

Tulis penjelasan terkait \tab~\ref{table:sample} di sini.

%\clearchapter
\input{src/01-body/06-kesimpulan}
%\clearchapter

%
% Daftar Pustaka
%%
% Daftar Pustaka
%
\chapter*{DAFTAR PUSTAKA}
\addChapter{DAFTAR PUSTAKA}
%
% Tambahkan pustaka yang digunakan setelah perintah berikut.
%
%\phantomsection %hack to add clickable section for pustaka
% \bibliographystyle{ieeetr}
% \bibliography{references}
\begin{otherlanguage}{english}
%\nocite{*}
% \subsection*{Buku}
% \noindent\textbf{Buku dan \textit{E-Book}}
%-----------------------------------------------------------------------------%
% \begingroup
\setstretch{1}
\setlength\bibitemsep{12pt}
\printbibliography[heading=none]
% \endgroup

\end{otherlanguage}

%\clearchapter

% Alternatif manajemen daftar pustaka dengan \bibliography
\bibliographystyle{apalike}
\bibliography{src/references}

%
% Lampiran
%
%\begin{appendix}
%	\input{etc/markLampiran}
%	%\clearchapter
%	\setcounter{page}{2}
%	%-----------------------------------------------------------------------------%
\begin{appendix}
% \renewcommand\pagenumbering[1]{} % no numbering
% \renewcommand*{\thepage}{\Roman}
    \input{etc/markLampiran}
%     \clearpage

    \pagestyle{lampiran}
%     \setcounter{page}{2}
%     \chapter{LAMPIRAN}

    \pagenumbering{Roman}
%     \pretocmd{\chapter}{%
%       \clearpage
%       \pagenumbering{Roman}%
%       \renewcommand*{\thepage}{\thechapter\arabic{page}}%
%     }{}{}

  \setlength{\leftskip}{0cm}
%----------------------------------------------------------------------------%

    \newlampiran{Surat Pengantar Penelitian}
%     \centerline{
%         \fbox{ %
%         \parbox{1.05\linewidth}{%
%           \includegraphics[width=\linewidth,trim={2cm 10.5cm 1.5cm 0.5cm},clip]{submission/lampiran_pengantar_penelitian.pdf}
%         %masukkan surat pengantar hasil scan dalam pdf
%           }
%         }
%         }
    \newpage

    % izin penelitian
    \phantomsection
    \newlampiran{Surat Izin Penelitian}
%     \centerline{
%         \frame{\includegraphics[width=1.05\linewidth]{submission/lampiran_surat_izin_penelitian.pdf}}
%         %masukkan surat izin penelitian dari lokasi
%         }
    \newpage


    % pedoman wawancara
    \phantomsection
    \newlampiran{Pedoman Wawancara}
    %
% Lampiran interview Guide
%
% @author  Slamet B. Aryo
% @version 1.10
%

% \newpage
% \pagenumbering{gobble} % enable only for proposal

% \noindent \bo{Lampiran Interview Guide}
\vspace{.5cm}

Berikut adalah pertanyaan wawancara terkait \judul.

\textbf{Kepada Pimpinan}

\begin{enumerate}
    \item Bagaimanakah sejarah berdirinya \lokus?
    \item Apa visi dan misi didirikannya \lokus?
\end{enumerate}

    \newpage

    % traskrip wawancara
    \newlampiran{Transkrip Wawancara}
    %
% Lampiran transkrip interview
%
% @author  Slamet B. Aryo
% @version 1.01
%
\begin{spacing}{1.125}
\noindent \underline{Bapak \narsumSatu~(Kepala \lokus)}

\vspace{1cm}

\begin{enumerate}
\setlength{\itemsep}{12pt}%
\item Bagaimana sejarah berdirinya \lokus?\\
\textit{Jawab:} \lipsum[41-49]

\end{enumerate}
\end{spacing}

    \newpage

    % dokumentasi
    \newlampiran{Dokumentasi}
    %\begin{center}
     \begin{figure}[H]
      \centering
        \includegraphics[width=0.8\textwidth]{assets/pics/makara.png}
        \caption*{Wawancara dengan \narsumSatu}
     \end{figure}

     \begin{figure}[H]
      \centering
        \includegraphics[width=0.8\textwidth]{assets/pics/creative_commons.png}
        \caption*{Wawancara dengan \narsumDua}
     \end{figure}
    %\end{center}
    \newpage


    % riwayat hidup
    \newlampiran{Daftar Riwayat Hidup}
    %
% Halaman riwayat hidup penulis
%
% @author  Slamet B. Aryo
% @version 1.01
%
{\centering \bo{RIWAYAT HIDUP PENULIS}\par}
\vspace*{1cm}

\begin{spacing}{1.125}
\noindent\bo{A. IDENTITAS}

\begin{enumerate}[labelindent=0em,labelsep=0em,itemsep=0em]
 \item \tabbox{Nama} : \penulis
 \item \tabbox{Tempat tanggal lahir} : Kota, Tgl Bulan Tahun
 \item \tabbox{Alamat rumah} : Alamat lengkapnya
 \item \tabbox{Alamat tinggal} : Alamat lengkapnya
 \item \tabbox{Nomor \textit{handphone}} : 08123xxxxxxx
 \item \tabbox{Email} : namakamu@mhs.uingusdur.ac.id
 \item \tabbox{Nama ayah} : Namanya
 \item \tabbox{Pekerjaan ayah} : Pekerjaan
 \item \tabbox{Nama ibu} : Namanya
 \item \tabbox{Pekerjaan ibu} : Pekerjaan
\end{enumerate}

\noindent\bo{B. RIWAYAT PENDIDIKAN}

\begin{enumerate}[labelindent=0em,labelsep=0em,itemsep=0em]
 \item \tabbox{SD} : SDN Namanya (20xx--20xx)
 \item \tabbox{SMP} : Nama SMP (20xx--20xx)
 \item \tabbox{SMA} : Nama SMA (20xx--20xx)
\end{enumerate}

\vspace*{1cm}
\begin{ttdkanan}
 Pekalongan, \tanggalPengesahan\\[4em]


 \penulis
\end{ttdkanan}


\end{spacing}

    \clearpage
\end{appendix}

%\end{appendix}

\end{document}
